\chapter{Axial Element}

\section{Stiffness matrix for an axial (bar, or truss) element}
\begin{fig}
  % Axial (bar/truss) element with end displacements and forces
\begin{tikzpicture}[scale=1.0,font=\small]
  \def\L{5}
  \draw[thick] (0,-0.15) rectangle (\L,0.15);
  \node[above] at (\L/2,0.15) {$EA$};
  \draw[->,thick] (-1.5,0) -- (-0.1,0);
  \node[above] at (-0.8,0.1) {$\delta_1\,(f_1)$};
  \draw[->,thick] (\L+0.1,0) -- (\L+1.5,0);
  \node[above] at (\L+0.8,0.1) {$\delta_2\,(f_2)$};
  \draw[dim] (0,-0.55) -- (\L,-0.55) node[midway,fill=white,inner sep=1pt] {$L$};
\end{tikzpicture}

\end{fig}
\begin{equation}
  \mat{f} = \kmat\bs{\delta} \quad\text{where}\quad
  \mat{f} = \Bmat{f_1\\ f_2},\quad \bs{\delta} = \Bmat{\delta_1\\ \delta_2}
  \quad\text{and}\quad \kmat = \frac{EA}{L}\bmat{1 & -1\\ -1 & 1}.
\end{equation}
\begin{equation*}
  \delta = \delta_2-\delta_1,\qquad f_2 = k\delta = \frac{EA}{L}(\delta_2-\delta_1)
  \quad\text{and}\quad f_1 = -f_2 = \frac{EA}{L}(\delta_1-\delta_2).
\end{equation*}

\section{Methods for deriving the element stiffness matrix \texorpdfstring{$\kmat$}{k}}
\begin{enumerate}[leftmargin=*]
  \item \textbf{Direct method}\\
    By definition, the $i$-th column of $\kmat$ is the $\mat{f}$ vector corresponding to $\delta_i = 1$, and other $\delta_j = 0$ ($j\neq i$).

    For example, take $\delta_1 = 1$, $\delta_2 = 0$; then the first column of the stiffness matrix can be obtained as shown.
    \begin{fig}
      % Direct method: impose delta_1 = 1, delta_2 = 0
\begin{tikzpicture}[scale=1.0,font=\small]
  \def\L{5}
  % undeformed bar
  \draw[thick] (0,0.35) rectangle (\L,0.65);
  % wall at the right end
  \fill[left color=black!15,right color=black!60] (\L,-0.5) rectangle (\L+0.3,1.0);
  % deformed bar (dotted), shifted 1 to the right at the left end
  \draw[densely dotted,thick] (0.6,-0.15) rectangle (\L,0.15);
  \draw[dim] (0,-0.15) -- (0.6,-0.15); \node[below] at (0.3,-0.2) {1};
  \draw (0,-0.3) -- (0,0.3);
  % end forces
  \draw[->,thick] (-2.0,0) -- (-0.1,0);
  \node[left] at (-2.1,0) {$k_{11}=\dfrac{EA}{L}$};
  \draw[->,thick] (\L+0.4,0) -- (\L+1.9,0);
  \node[right] at (\L+2.0,0) {$k_{21}=-\dfrac{EA}{L}$};
\end{tikzpicture}

    \end{fig}
  \item \textbf{Energy method}\\
    The strain energy $U$ in terms of $\bs{\delta}$ for a linearly elastic member is a quadratic form. Hence
    \begin{equation}
      U = \tfrac12\,\bs{\delta}^\tr\kmat\bs{\delta}.
    \end{equation}
    This is equivalent to
    \begin{equation}
      k_{ij} = \frac{\partial^2 U}{\partial\delta_i\,\partial\delta_j}.
    \end{equation}
    For a truss element, $U$ can be easily shown to be
    \begin{equation}
      U = \frac12\frac{EA}{L}(\Delta\delta)^2 = \frac12\frac{EA}{L}(\delta_2-\delta_1)^2.
    \end{equation}
  \item \textbf{Principle of virtual displacement}, or finite element approach employing shape functions (we shall discuss this method in Ch.~\ref{ch:pvd2}).
\end{enumerate}

\section{Global stiffness matrix for an axial element}
\emph{Stiffness matrix ready for assembling.}
\begin{equation}
  \Bmat{F_{x1}\\ F_{y1}\\ F_{x2}\\ F_{y2}} = \frac{EA}{L}
  \bmat{c^2 & sc & -c^2 & -sc\\ sc & s^2 & -sc & -s^2\\ -c^2 & -sc & c^2 & sc\\ -sc & -s^2 & sc & s^2}
  \Bmat{u_1\\ v_1\\ u_2\\ v_2}
  \label{eq:truss-global-k}
\end{equation}
\begin{minipage}{0.68\textwidth}
where $c = \cos\phi$, $s = \sin\phi$, and $\phi$ is the angle between the local ($x$) and the global ($X$) axes, which is positive if the $x$ axis is rotated counterclockwise with respect to the $X$ axis. Methods for deriving the local stiffness matrix can also be used to obtain the above global stiffness matrix.
\end{minipage}\hfill
\begin{minipage}{0.28\textwidth}
  \centering
  % Local (x,y) axes rotated by phi from global (X,Y) axes
\begin{tikzpicture}[scale=1.0,font=\small]
  \draw[->,thick] (0,0) -- (3.2,0) node[right] {$X$};
  \draw[->,thick] (0,0) -- (0,3.2) node[above] {$Y$};
  \coordinate (O) at (1.2,0.7);
  \draw[dashed] (O) -- ++(1.6,0);
  \draw[->,thick] (O) -- ++(35:2.0) node[above right] {$x$};
  \draw[->,thick] (O) -- ++(125:1.4) node[above left] {$y$};
  \draw[->] ($(O)+(1.1,0)$) arc (0:35:1.1);
  \node at ($(O)+(1.35,0.35)$) {$\phi$};
\end{tikzpicture}

\end{minipage}

Note that the strain energy is
\begin{equation}
  U = \frac12\frac{EA}{L}(\delta_2-\delta_1)^2
    = \frac12\frac{EA}{L}\langle\delta_1\;\;\delta_2\rangle\bmat{1 & -1\\ -1 & 1}\Bmat{\delta_1\\ \delta_2}.
\end{equation}
The member displacements in the local axes, $\delta_1$ and $\delta_2$, and the member displacements in the global axes, $u_1$, $v_1$, $u_2$ and $v_2$, are shown in the figure below.
\begin{fig}
  % Member displacements in local (delta) and global (u,v) axes
\begin{tikzpicture}[scale=1.0,font=\small]
  \coordinate (A) at (0,0); \coordinate (B) at (35:4.2);
  \draw[member] (A) -- (B);
  \draw[dashed] (A) -- ++(215:0.9); \draw[dashed] (B) -- ++(35:0.9);
  % node 1
  \draw[->,thick] (A) -- ++(1.4,0) node[below right] {$u_1$};
  \draw[->,thick] (A) -- ++(0,1.4) node[above] {$v_1$};
  \draw[->,thick] ($(A)+(215:0.9)$) -- ($(A)+(215:0.05)$);
  \node[left] at ($(A)+(215:0.7)$) {$\delta_1$};
  \draw[->] ($(A)+(0.9,0)$) arc (0:35:0.9); \node at ($(A)+(1.15,0.32)$) {$\phi$};
  % node 2
  \draw[->,thick] (B) -- ++(1.4,0) node[right] {$u_2$};
  \draw[->,thick] (B) -- ++(0,1.4) node[above] {$v_2$};
  \draw[->,thick] (B) -- ++(35:1.0) node[above right] {$\delta_2$};
\end{tikzpicture}

\end{fig}

Relationship between the displacements in the local axes and the displacements in the global axes:
\begin{equation}
  \bs{\delta} = \Bmat{\delta_1\\ \delta_2} = \bmat{c & s & 0 & 0\\ 0 & 0 & c & s}\Bmat{u_1\\ v_1\\ u_2\\ v_2} = \Tmat\bs{\Delta}.
\end{equation}
Writing the strain energy $U$ in terms of $\bs{\delta}$ and $\bs{\Delta}$:
\begin{equation}
  U = \frac12\bs{\delta}^\tr\kmat_l\bs{\delta} = \frac12\bs{\Delta}^\tr\Tmat^\tr\kmat_l\Tmat\bs{\Delta}
    = \frac12\bs{\Delta}^\tr\kmat_g\bs{\Delta},
  \qquad\text{where } \kmat_g = \Tmat^\tr\kmat_l\Tmat.
\end{equation}
Therefore,
\begin{equation}
  U = \frac12\Bmat{u_1\\ v_1\\ u_2\\ v_2}^{\!\tr}
  \bmat{c & 0\\ s & 0\\ 0 & c\\ 0 & s}\frac{EA}{L}\bmat{1 & -1\\ -1 & 1}\bmat{c & s & 0 & 0\\ 0 & 0 & c & s}
  \Bmat{u_1\\ v_1\\ u_2\\ v_2}
  = \frac12\Bmat{u_1\\ v_1\\ u_2\\ v_2}^{\!\tr}[k_g]\Bmat{u_1\\ v_1\\ u_2\\ v_2},
\end{equation}
where $[k_g]$ is the global stiffness matrix given above.

\section{General features of an element stiffness matrix}
\begin{enumerate}[leftmargin=*]
  \item \textbf{Symmetric}
  \item \textbf{Singular}\\
    Number of zero eigenvalues $=$ number of rigid-body modes.\\
    E.g.\ for the local stiffness matrix of an axial element, the number is 1 (one translation), and for the global one, the number is 3 (two translations plus one rotation). For example, let's take a look at the rigid-body rotation (w.r.t.\ node 1) mode, $\bs{\Delta}_r$, where
    \begin{fig}
      \begin{minipage}{0.35\textwidth}
        \begin{equation*}
          \bs{\Delta}_r = \Bmat{0\\ 0\\ -L\theta_r\sin\phi\\ L\theta_r\cos\phi}
        \end{equation*}
      \end{minipage}
      \begin{minipage}{0.55\textwidth}
        \centering% Rigid-body rotation mode about node 1
\begin{tikzpicture}[scale=1.0,font=\small]
  \coordinate (A) at (0,0); \coordinate (B) at (30:4.2); \coordinate (Br) at (55:4.2);
  \draw[member] (A) -- (B); \draw[member] (A) -- (Br);
  \draw[dashed] (B) -- (Br);
  \node[left] at ($(B)!0.5!(Br)+(-0.05,0)$) {$L\theta_r$};
  \draw[->,thick] (A) -- ++(1.4,0) node[below right] {$u_1$};
  \draw[->,thick] (A) -- ++(0,1.4) node[left] {$v_1$};
  \draw[->,thick] (B) -- ++(1.2,0) node[right] {$u_2$};
  \draw[->,thick] (B) -- ++(0,1.2) node[above] {$v_2$};
  \draw[->] ($(A)+(1.0,0)$) arc (0:30:1.0); \node at ($(A)+(1.25,0.28)$) {$\phi$};
  \draw[->] ($(A)+(30:1.7)$) arc (30:55:1.7); \node at ($(A)+(42:2.0)$) {$\theta_r$};
\end{tikzpicture}

      \end{minipage}
    \end{fig}
  \item Each column of the element stiffness matrix constitutes a set of member end forces that are in \textbf{equilibrium}. Therefore, $\F = \kmat\bs{\Delta}$ is also a set of forces in equilibrium.

    For example, $v_2 = 1$ and others equal 0. The axial force is equal to $(EA/L)\sin\phi$. With the geometric relationship, we get
    \begin{equation}
      \Bmat{F_{x1}\\ F_{y1}\\ F_{x2}\\ F_{y2}} = \frac{EA}{L}\Bmat{-sc\\ -s^2\\ sc\\ s^2}.
    \end{equation}
    This vector is the 4th column of the element stiffness matrix.
    \begin{fig}
      % Fourth column of the global stiffness matrix: v2 = 1, others 0
\begin{tikzpicture}[scale=1.0,font=\small]
  % --- left: geometry ---
  \begin{scope}
    \coordinate (A) at (0,0); \coordinate (B) at (25:4.0);
    \draw[member] (A) -- (B);
    \draw[dashed,thick] (A) -- ($(B)+(0,1.0)$);
    \draw[->,thick] (B) -- ($(B)+(0,1.0)$);
    \node[left] at ($(B)+(-0.05,0.55)$) {1};
    % component of the unit displacement along the member
    \draw[densely dotted] ($(B)+(0,1.0)$) -- ($(B)+(25:{sin(25)})$);
    \draw[dim] ($(B)+(-0.1,0)$) -- ($(B)+(25:{sin(25)})+(-0.1,0)$);
    \node[below right,inner sep=1pt] at ($(B)+(0.05,-0.05)$) {\scriptsize $1\cdot\sin\phi$};
    \draw[->] ($(A)+(1.4,0)$) arc (0:25:1.4); \node at ($(A)+(1.7,0.3)$) {$\phi$};
    \draw (A) -- ++(1.8,0);
  \end{scope}
  % --- right: end forces ---
  \begin{scope}[xshift=7.5cm]
    \coordinate (A) at (0,0); \coordinate (B) at (25:4.0);
    \draw[member] (A) -- (B);
    \draw[->,thick] (A) -- ++(0,1.3) node[right] {$-\dfrac{EA}{L}\cdot s^2$};
    \draw[->,thick] (A) -- ++(1.2,0) node[below] {$-\dfrac{EA}{L}\cdot sc$};
    \draw[->,thick] (B) -- ++(0,1.3) node[right] {$\dfrac{EA}{L}\cdot s^2$};
    \draw[->,thick] (B) -- ++(1.3,0) node[below right] {$\dfrac{EA}{L}\cdot sc$};
    \draw[->] ($(A)+(1.0,0)$) arc (0:25:1.0); \node at ($(A)+(1.3,0.35)$) {$\phi$};
  \end{scope}
\end{tikzpicture}

    \end{fig}
\end{enumerate}
